{\small\textbf{Across borders (Zheng, 1:00) [20:00–21:00]}}\par\smallskip
Second, borders.
\begin{itemize}\setlength\itemsep{0pt}
\item \textcolor{navy}{\textbf{[def]}} Under EU reciprocity, any foreign bank lending into the Netherlands must also hold the Dutch two percent on those loans.
\item That limits leakage, where lending simply moves to foreign banks, a known problem with national capital rules. Lending can still leak to non-banks, though.
\end{itemize}
\par\smallskip
And the Netherlands is part of a wider shift: Germany, France, Ireland and Belgium raised their buffers in 2023–24, and Portugal and Spain follow in 2026. In 2023 the Netherlands had the highest neutral rate in the banking union. The ECB's own surveys describe the same shift: authorities now want to build releasable capital early, rather than wait for a credit boom that the indicators may never show.
